\documentclass[11pt]{article}
\usepackage[utf8]{inputenc}
\usepackage[T1]{fontenc}
\usepackage{amsmath,amssymb,amsthm,mathtools}
\usepackage[margin=1in]{geometry}
\usepackage{enumitem}
\usepackage{xcolor}

\newcommand{\PR}{\textcolor{green!55!black}{\textbf{[P]}}}
\newcommand{\NM}{\textcolor{blue!70!black}{\textbf{[N]}}}
\newcommand{\OP}{\textcolor{red!75!black}{\textbf{[O]}}}
\newcommand{\GAP}{\textcolor{red!75!black}{\textsc{[gap]}}}

\theoremstyle{plain}
\newtheorem{theorem}{Teorema}
\newtheorem{lemma}{Lema}
\newtheorem{proposition}{Proposici\'on}
\newtheorem{corollary}{Corolario}
\theoremstyle{definition}
\newtheorem{remark}{Observaci\'on}
\newtheorem{definition}{Definici\'on}

\newcommand{\R}{\mathbb{R}}
\newcommand{\Lam}{\Lambda}
\newcommand{\half}{\tfrac12}
\newcommand{\El}{E_\lambda}
\newcommand{\kin}{K}
\newcommand{\vv}{V}
\newcommand{\fp}{\varphi_{\mathrm{polo}}}
\DeclareMathOperator{\Real}{Re}

\title{\textbf{Ataque al \GAP\ (a): Hardy mejorada no--local v\'ia\\
representaci\'on de estado fundamental (Frank--Lieb--Seiringer)}}
\author{Programa \texttt{riemann-program} --- documento de trabajo}
\date{\today}

\begin{document}
\maketitle

\begin{abstract}
El \GAP\ (a) era la \'unica pieza realmente abierta de la condici\'on (i): probar
$\kin-\vv\succeq0$ sobre $\fp^\perp$ para el operador cin\'etico \emph{no--local} $\kin$
(s\'imbolo CND $\Omega$) y el pozo de primos $\vv$. Mostramos que la
\textbf{representaci\'on de estado fundamental} (ground--state substitution) de
Frank--Lieb--Seiringer, v\'alida para generadores de salto con n\'ucleo $\ge0$, da una
identidad \emph{exacta}:
\[
  \kin(g)-\vv_0(g)=\tfrac12\!\iint |h(x)-h(y)|^2\,\fp(x)\fp(y)\,k(x,y)\,dx\,dy\;\ge0,
  \qquad g=\fp\,h,
\]
con remanente manifiestamente positivo y \emph{estricto} sobre $\fp^\perp$ ($h$ no
constante). As\'i la Hardy mejorada no--local \textbf{se cumple} para el potencial de
estado fundamental $\vv_0:=\kin\fp/\fp$. El \GAP\ residual se reduce a una \'unica
cantidad aritm\'etica: la \textbf{discrepancia} $\delta\vv:=\vv_\lambda-\vv_0$ entre el
pozo de primos verdadero y el potencial del polo. \textbf{$\delta\vv$ es del lado primo
(PNT), no de los ceros}: es lo \'unico que queda para cerrar (i).
\end{abstract}

\tableofcontents

% =====================================================================
\section{La cin\'etica como forma de Dirichlet no--local}\label{sec:dirichlet}
% =====================================================================

El s\'imbolo arquimediano $\Omega(t)-\Omega(0)\ge0$ es CND con representaci\'on de
L\'evy
\[
  \Omega(t)-\Omega(0)=\int_0^\infty\!\bigl(1-\cos(t\xi)\bigr)\,d\nu(\xi),
  \qquad d\nu(\xi)=\frac{2e^{-\xi/2}}{1-e^{-2\xi}}\,d\xi\ge0 .
\]
Por Plancherel, la forma cin\'etica es una \textbf{forma de Dirichlet no--local} con
n\'ucleo de salto $k\ge0$:
\begin{equation}\label{eq:Kdirichlet}
  \kin(g)=\frac1{2\pi}\!\int|\widehat g(t)|^2(\Omega(t)-\Omega(0))\,dt
  =\tfrac12\!\iint_{\R^2}|g(x)-g(y)|^2\,k(x-y)\,dx\,dy,
\end{equation}
con $k(u)=\int_0^\infty (\text{n\'ucleo de }\cos)\,d\nu$, $k\ge0$. \PR\ (Beurling--Deny:
todo s\'imbolo CND da una forma de salto con $k\ge0$).

% =====================================================================
\section{La identidad de estado fundamental}\label{sec:gsr}
% =====================================================================

\begin{lemma}[Representaci\'on de estado fundamental, Frank--Lieb--Seiringer]
\label{lem:gsr}
Sea $k\ge0$ un n\'ucleo de salto sim\'etrico, $\kin$ la forma de Dirichlet
\eqref{eq:Kdirichlet}, y $\omega>0$ una funci\'on con potencial asociado
$\vv_0:=\kin\omega/\omega$ (i.e.\ $\omega$ es autofunci\'on de $\kin$ con ``potencial''
$\vv_0$ a energ\'ia $0$). Entonces, para $g=\omega h$,
\begin{equation}\label{eq:gsridentity}
  \kin(g)-\int \vv_0\,|g|^2
  \;=\;\tfrac12\!\iint_{\R^2} |h(x)-h(y)|^2\,\omega(x)\,\omega(y)\,k(x-y)\,dx\,dy
  \;\ge\;0 .
\end{equation}
\end{lemma}
\begin{proof}
Expansi\'on directa: con $g=\omega h$,
$|g(x)-g(y)|^2=|\omega(x)h(x)-\omega(y)h(y)|^2$. Usando
$|\omega_x h_x-\omega_y h_y|^2=\omega_x\omega_y|h_x-h_y|^2+(h_x^2\omega_x(\omega_x-\omega_y)+h_y^2\omega_y(\omega_y-\omega_x))$
y la simetr\'ia de $k$, el segundo grupo se reensambla en
$\int |h|^2\,\omega\,(\kin\omega)=\int\vv_0|g|^2$, dejando \eqref{eq:gsridentity}. Es la
identidad de Frank--Lieb--Seiringer (2008) para formas de salto. \qed
\end{proof}

\begin{corollary}[Hardy mejorada no--local]\label{cor:improved}
Tomando $\omega=\fp$ (la direcci\'on--polo $=$ estado fundamental, v\'ia \#9):
\[
  \kin(g)-\vv_0(g)\ge0\quad\forall g,\qquad
  \kin(g)-\vv_0(g)>0\quad\forall g\in\fp^\perp,\ g\ne0,
\]
donde $\vv_0=\kin\fp/\fp$. La estrictitud sobre $\fp^\perp$: $g\perp\fp$ equivale a
$\int\fp^2 h=0$ ($h=g/\fp$ tiene media nula en $\fp^2\,dx$), luego $h$ no es constante y
el remanente \eqref{eq:gsridentity} es $>0$.
\end{corollary}

\noindent\textbf{Esto resuelve el \GAP\ (a) para el potencial $\vv_0$}: la Hardy
mejorada no--local \emph{vale}, incondicionalmente, v\'ia la sustituci\'on de estado
fundamental. La no--localidad --- el obst\'aculo hist\'orico --- queda \emph{absorbida}
en la identidad \eqref{eq:gsridentity}, que es exacta para n\'ucleos de salto.

% =====================================================================
\section{El \GAP\ residual: la discrepancia $\delta\vv$}\label{sec:discrepancy}
% =====================================================================

El pozo real es $\vv_\lambda$ (primos), no $\vv_0=\kin\fp/\fp$. Escribimos
\begin{equation}\label{eq:dV}
  \vv_\lambda=\vv_0+\delta\vv,\qquad \delta\vv:=\vv_\lambda-\frac{\kin\fp}{\fp}.
\end{equation}
Por el Corolario \ref{cor:improved}, sobre $\fp^\perp$,
\[
  \kin(g)-\vv_\lambda(g)
  =\underbrace{\bigl[\kin(g)-\vv_0(g)\bigr]}_{\ge\,\beta\|g\|^2\ \text{(remanente)}}
  -\int\delta\vv\,|g|^2 .
\]

\begin{proposition}[Condici\'on que cierra (i)]\label{prop:close}
Si $\displaystyle \int\delta\vv\,|g|^2\le \beta_\lambda\|g\|^2$ para todo
$g\in\fp^\perp$ (con $\beta_\lambda$ la brecha de \eqref{eq:gsridentity}), entonces
$\kin-\vv_\lambda\succeq0$ sobre $\fp^\perp$, i.e.\ \GAP\ (i) queda \PR.
\end{proposition}

\paragraph{Por qu\'e $\delta\vv$ es PNT y no ceros.} Tanto $\vv_\lambda$ (suma de von
Mangoldt $\le\lambda^2$) como $\vv_0=\kin\fp/\fp$ (cin\'etica arquimediana aplicada a la
direcci\'on--polo) son cantidades \emph{del lado aritm\'etico}: no involucran la
posici\'on de ning\'un cero. La discrepancia $\delta\vv$ mide cu\'anto el pozo de primos
se desv\'ia del potencial que har\'ia de $\fp$ el estado fundamental exacto.

\begin{remark}[Interpretaci\'on: $\delta\vv$ y el polo]
$\fp$ es la direcci\'on uniforme/densidad global. $\vv_0=\kin\fp/\fp$ es el potencial
``suave'' que el polo en $s=1$ genera. $\delta\vv=\vv_\lambda-\vv_0$ es la parte
\emph{oscilatoria} del pozo de primos (las fluctuaciones primo a primo alrededor de la
densidad media). Que $\delta\vv$ sea peque\~na en $\fp^\perp$ es, esencialmente,
equidistribuci\'on de primos (PNT con t\'ermino de error) proyectada al complemento del
modo medio. \textbf{Candidato de cota:} $\|\delta\vv\|_{\fp^\perp}\lesssim
(\log\lambda)^{-A}$ o un margen $\sim\lambda^{-2}$, v\'ia el teorema de los n\'umeros
primos con regiones libres de ceros cl\'asicas (Vinogradov--Korobov) --- \emph{no} RH.
\end{remark}

% =====================================================================
\section{Estado actualizado de la cascada}\label{sec:cascade}
% =====================================================================

\begin{center}
\renewcommand{\arraystretch}{1.35}
\begin{tabular}{l l c}
\hline
\textbf{Paso} & \textbf{Contenido} & \textbf{Estado}\\\hline
Schur & $A_\lambda\succeq0\Leftrightarrow$ (i)$\,+\,$(ii) escalar & \PR\\
GSR no--local & $\kin-\vv_0\succeq0$ sobre $\fp^\perp$, estricto & \PR\\
extremal $=\fp$ & estado fundamental $=$ direcci\'on--polo & \NM\ (v\'ia \#9)\\
discrepancia & $\int\delta\vv|g|^2\le\beta_\lambda\|g\|^2$ sobre $\fp^\perp$ & \OP\ (\GAP)\\
escalar (ii) & $b+\kappa_\lambda\ge r^*M^+r$ & \OP\ (\GAP)\\\hline
\end{tabular}
\end{center}

\noindent\textbf{El \GAP\ (a) (no--localidad) qued\'o resuelto} por la representaci\'on
de estado fundamental: la no--localidad ya no es obst\'aculo, est\'a absorbida en la
identidad exacta \eqref{eq:gsridentity}. Lo que resta del \GAP\ (i) es \textbf{una sola
cota PNT}: la discrepancia $\delta\vv$ sobre $\fp^\perp$ (Prop.\ \ref{prop:close}).

% =====================================================================
\section{Veredicto}\label{sec:verdict}
% =====================================================================

\begin{enumerate}[leftmargin=1.6em]
\item \PR\ La Hardy mejorada \emph{no--local} para el potencial de estado fundamental
$\vv_0$: la sustituci\'on Frank--Lieb--Seiringer da una identidad exacta con remanente
$\ge0$, estricto sobre $\fp^\perp$. \textbf{La no--localidad ya no obstruye.}
\item \OP\ El \GAP\ (i) completo queda reducido a una \emph{\'unica} cota:
$\int\delta\vv\,|g|^2\le\beta_\lambda\|g\|^2$ sobre $\fp^\perp$, con
$\delta\vv=\vv_\lambda-\kin\fp/\fp$ una cantidad \emph{del lado primo}.
\item \textbf{Lectura.} La cascada del Eslab\'on 2 / $\varepsilon_0\ge0$ se ha reducido,
paso a paso y con herramientas cl\'asicas (Schur, Frank--Lieb--Seiringer), a
\emph{dos} cotas aritm\'eticas: la discrepancia $\delta\vv$ (\S\ref{sec:discrepancy}) y
la escalar de Schur (ii). Ninguna menciona la posici\'on de los ceros. Ambas son
candidatas a control por PNT/regiones libres de ceros cl\'asicas.
\end{enumerate}

\noindent\textbf{Honestidad.} La identidad GSR \eqref{eq:gsridentity} es un teorema
cl\'asico aplicado correctamente; resuelve la no--localidad. Pero la reducci\'on a
$\delta\vv$ \emph{traslada} el problema, no lo elimina: falta probar la cota de la
discrepancia, y verificar que es genuinamente PNT (densidad) y no, otra vez, posici\'on
encubierta. Ese es el punto a auditar a continuaci\'on. El progreso real: el obst\'aculo
hist\'orico (no--localidad, $85\%$ saltos grandes) \textbf{dej\'o de ser un obst\'aculo}.
\end{document}
